
\section{Introduction}
The increasing deployment of Ethernet-based infrastructures has weakened the assumption that the data-link layer is inherently trustworthy. In these settings, link-layer traffic is exposed to adversarial conditions, and cryptographic protection is needed to preserve confidentiality, integrity, and replay resistance. The IEEE 802.1AE MACsec standard~\cite{ieee8021ae} addresses this need by defining hop-by-hop protection for Ethernet frames, while the MACsec Key Agreement protocol (MKA)~\cite{ieee8021x} manages the establishment and maintenance of the cryptographic state used by MACsec.

This paper addresses that gap by presenting a formal analysis of a simplified two-party MKA exchange using the Tamarin prover. Our model captures the initial session establishment and a subsequent rekey round, and we verify essential security properties under a standard Dolev--Yao adversary model. The analysis confirms the expected guarantees when the CAK remains uncompromised, while also revealing a structural weakness in the protocol's acceptance logic.

\paragraph{Contributions.} This work makes three contributions. First, it constructs a Tamarin model of a two-party MKA exchange that captures the initial session and a rekey round, including the monotonic MN/KN counters. Second, it formalizes and verifies core security properties such as executability, agreement, authentication, secrecy, ordering, and freshness. Third, it identifies and demonstrates a structural weakness in the protocol's acceptance logic: a malicious or compromised Key Server that possesses the CAK can inject an arbitrary SAK that the Server accepts.

\subsection{Related Work}
\label{sec:related}

Formal verification has become standard for protocol analysis. Classical examples include Lowe's analysis of Needham--Schroeder with FDR~\cite{lowe1996breaking}, Blanchet's ProVerif-based work on authentication protocols~\cite{blanchet2001efficient}, and later Tamarin studies of TLS~1.3~\cite{cremers2017comprehensive}, Signal~\cite{kobeissi2017automated}, and WPA2~\cite{vanhoef2017key}. Within the IEEE 802 family, the closest related work concerns WPA2's 4-way handshake; however, the MKA key-agreement logic has received comparatively little machine-checked analysis.

Our contribution is therefore complementary. To the best of our knowledge, this is the first machine-verified symbolic analysis of a simplified MKA exchange. The closest methodological precedent is the Tamarin analysis of TLS~1.3 by Cremers et al.~\cite{cremers2017comprehensive}, which similarly verifies authentication, secrecy, and ordering properties. Our work applies the same general methodology to the Ethernet security domain, where the role of the Key Server introduces a distinct acceptance semantics that is not present in classical protocol examples.

\paragraph{Organization.} Section~\ref{sec:background} provides an overview of MKA and its security requirements. Section~\ref{sec:analysis} introduces the Tamarin Prover framework. Section~\ref{sec:model} presents our formal Tamarin model and details the fourteen security lemmas. Finally, Section~\ref{sec:conclusion} concludes the paper by synthesizing the verification results, structural trust limitations, and future research directions.


